\documentclass[11pt]{article}
\usepackage[T1]{fontenc}
\usepackage[utf8]{inputenc}
\usepackage{lmodern}
\usepackage[letterpaper,margin=1in,headheight=14pt]{geometry}
\usepackage{amsmath,amssymb,amsthm,mathtools,mathrsfs}
\usepackage{booktabs,longtable,array,enumitem}
\usepackage{microtype}
\usepackage[dvipsnames]{xcolor}
\usepackage{fancyhdr}
\usepackage{hyperref}
\definecolor{ink}{HTML}{17334A}
\definecolor{accent}{HTML}{176B71}
\hypersetup{colorlinks=true,linkcolor=accent,citecolor=accent,urlcolor=accent,
 pdftitle={Direct quadratic packing for a ten-integer equation equivalent to ZFC inconsistency},
 pdfauthor={AI-assisted research prepared for the ProveIt project}}
\pagestyle{fancy}
\fancyhf{}
\fancyhead[L]{\small Direct quadratic packing for \(\operatorname{Con}(\mathrm{ZFC})\)}
\fancyhead[R]{\small 23 October 2026}
\fancyfoot[C]{\thepage}
\renewcommand{\headrulewidth}{0.3pt}
\setlength{\parskip}{4pt}
\setlength{\parindent}{0pt}
\setlist{itemsep=3pt,topsep=5pt}
\allowdisplaybreaks[2]
\emergencystretch=2em
\newcommand{\Z}{\mathbb Z}
\newcommand{\N}{\mathbb N}
\newcommand{\Con}{\operatorname{Con}}
\newcommand{\Sq}{\operatorname{Sq}}
\newcommand{\pc}{\operatorname{pop}}
\newcommand{\BB}{\mathcal B}
\newcommand{\capC}{\mathscr C}
\newcommand{\Dcode}{D_{\mathrm c}}
\newcommand{\ZFC}{\mathrm{ZFC}}
\newcommand{\file}[1]{\nolinkurl{#1}}
\newtheorem{theorem}{Theorem}[section]
\newtheorem{proposition}[theorem]{Proposition}
\newtheorem{lemma}[theorem]{Lemma}
\newtheorem{corollary}[theorem]{Corollary}
\theoremstyle{definition}
\newtheorem{definition}[theorem]{Definition}
\theoremstyle{remark}
\newtheorem{remark}[theorem]{Remark}
\numberwithin{equation}{section}

\begin{document}
\begin{center}
{\small\scshape ProveIt research contribution}\\[8pt]
{\LARGE\bfseries\color{ink} Direct quadratic packing for a\\[4pt]
ten-integer consistency equation}\\[10pt]
{\large A polynomial of exact degree \(237\,520\,456\) in 782 operations}\\[10pt]
AI-assisted research prepared for Vladimir Reshetnikov\\
23 October 2026
\end{center}

\begin{abstract}
We replace the source coding in the repository's ten-integer Diophantine
representation of ZFC inconsistency. The replacement encodes the 22 quadratic
parent equations directly, using 32 Sidon positions and 22 independently
tested target digits. It never forms the quartic sum of squares before
packing. The resulting completely fixed polynomial
\[
 F_{10}(f,g,h,k,l,w,x,y,z,m)\in\Z[f,g,h,k,l,w,x,y,z,m]
\]
has exact degree \(237\,520\,456\) and a division-free circuit with 782 binary
arithmetic operations. On the inherited, unchanged K215/compiler coefficient
tuple, ordinary mathematics formalizable in ZFC proves
\[
 \bigl(\nexists\,\mathbf v\in\Z^{10}:F_{10}(\mathbf v)=0\bigr)
 \quad\longleftrightarrow\quad \Con(\ZFC).
\]
Compared with package 31 at the pinned repository revision, the number and
domain of unknowns are unchanged, the degree falls by a factor exceeding
\(2.47\cdot10^{18}\), and the circuit loses 660 operations.

The new ingredient is a proved source reduction for quadratic systems.
Sun's Lucas and binomial bridges, including his 2026 preprint, and the
repository's modified three-square finalizer are retained with explicit
attribution. The archive contains the polynomial circuit, all fixed
coefficient recipes and payloads, the literal machine, exact algebraic
audits, independent circuit checks, and this article. This is not an
end-to-end proof-assistant certificate or a claim of global minimality.
\end{abstract}

\vspace{4pt}
\begin{center}
\small
\begin{tabular}{@{}lrr@{}}
\toprule
Metric & Package 31 & This construction\\
\midrule
Unrestricted integer unknowns & 10 & 10\\
Exact degree & \(126496\cdot5^{31}+27064\) & \(237\,520\,456\)\\
Binary \(+,-,\times\) operations & 1,442 & 782\\
Native coefficient tuple & \(121\times8801\) & unchanged\\
\bottomrule
\end{tabular}
\end{center}

\newpage
{\small\setlength{\parskip}{0pt}\tableofcontents}
\newpage

\section{The result and the precise comparison}
\label{sec:result}

The reference tree is Vladimir Reshetnikov's ProveIt repository at commit
\[
\texttt{22e75dd6b87b1b199bb08c778d7cf075ae5aae08}.
\]
The comparison is with the current collection and, specifically, its
\file{31-sun10} package, rather than an earlier superseded wave
\cite{repo,package31}. The copied collection description is
\file{upstream/COMPARISON_AT_PIN.md}.

Fix the repository's effective proof calculus for ZFC. In this article
\[
 \Con(\ZFC)\;:\Longleftrightarrow\;
 \neg\exists p\;\operatorname{Proof}_{\ZFC}(p,\bot)
\]
means absence of a finite refutation in that calculus. This is the sentence
calibrated by the endpoint
\file{BProv ZFCax_s [] Form.fBot} in \file{KEndpoint215.lean}.
We use \(\N=\{0,1,2,\ldots\}\); positive integers are denoted by
\(\Z_{>0}\).

\begin{theorem}[Main construction]\label{thm:main}
The definitions in Sections \ref{sec:parent}--\ref{sec:equation}, with the
fixed data supplied in the archive, determine one polynomial \(F_{10}\)
over the integers in exactly the ten unknowns
\[
 (f,g,h,k,l,w,x,y,z,m)\in\Z^{10}.
\]
Its exact total degree is \(237\,520\,456\). The supplied circuit has
221 additions, 25 subtractions, and 536 multiplications.
Composing the new reduction with the inherited K215/native-compiler
equivalence and the cited Lucas theorems gives a proof formalizable in ZFC of
\begin{equation}\label{eq:main}
 \exists\mathbf v\in\Z^{10}\;F_{10}(\mathbf v)=0
 \quad\longleftrightarrow\quad
 \exists p\;\operatorname{Proof}_{\ZFC}(p,\bot).
\end{equation}
Consequently ZFC proves that absence of an integer zero is equivalent to
\(\Con(\ZFC)\).
\end{theorem}

All the auxiliary expressions below are polynomial abbreviations. They
introduce no further existential unknowns. All exponents, finite sum limits,
and coefficients are fixed independently of whether ZFC has a refutation.
In particular, no unspecified recursively enumerable index or
consistency-dependent choice occurs in the definition.

\subsection{What improves, and what is retained}

The exact old degree is
\[
126496\cdot5^{31}+27064
=589042901992797851562527064.
\]
The ratio of this degree to the new degree is approximately
\(2.4799670391\cdot10^{18}\). The operation reduction is
\(660/1442\), approximately \(45.77\%\).

These are comparisons on the \emph{same ten-unrestricted-integer branch}.
The repository also contains much smaller operation counts with more
unknowns or with positivity restrictions, low-degree systems with many
witnesses, and nine-positive-variable forms \cite{repo}. The present
result does not dominate all of these measures. The ten-variable
unrestricted-integer reduction itself is Sun's result \cite{sun2026}.
Package 31's quintic divisibility merger, coprime square merger, and
structural separating weights are also retained, not claimed anew.

The new synthesis changes the input to that finalizer:
\[
\begin{array}{c@{\quad}c@{\quad}c}
\text{old: quadratic rows}
 &\longrightarrow\text{one quartic sum of squares}
 &\longrightarrow\text{positions }5^j,\\[3pt]
\text{new: quadratic rows}
 &\longrightarrow\text{Sidon positions}
 &\longrightarrow\text{all target digits at once}.
\end{array}
\]
The first row loses the low degree of the system before coding it.
The second row preserves it and uses the same packed witness for every
equation.

\subsection{The mathematical dependency boundary}

The proof has three distinct layers.
First, the inherited construction relates the literal K215 machine,
its marked Wang language, a native-window/Pell compiler, and the positive
quadratic parent system. Second, Sections \ref{sec:packing} and
\ref{sec:source} prove the new simultaneous-digit source reduction.
Third, Sections \ref{sec:equation} and \ref{sec:finalproof} compose that
source with Sun's bridges and the retained finalizer.

The repository provides formal K215 endpoint sources. This article does
not run their full proof-assistant dependency tree or infer ZFC provability
from a blanket assertion that ZFC proves a proof assistant sound.
The mathematical route uses the substantive inherited K215 calibration,
formalizable in ZFC. Besides finite simulation and proof enumeration,
that calibration includes a seed/quotient model construction,
first-order completeness, and the ZF--ZFC refutability equivalence
obtained through the constructible universe. The finite programs in this archive verify
transcriptions and algebraic claims; they do not prove those entire
inherited semantic theorems or establish \(\Con(\ZFC)\).

\section{The fixed quadratic parent}
\label{sec:parent}

We retain the complete coefficient tuple used by package 31. Its four
positive integer ports are
\[
 (\beta,K_0,M_C,M_F),
\]
where \(M_F\) is the \emph{already shifted} high mask.
Appendix \ref{app:constants} gives the finite recipe; the archive also
materializes all four integers. Their respective bit lengths are
\[
 (1064921,\;1070004,\;1064921,\;1064922).
\]
Keeping the whole tuple unchanged avoids any transfer claim between
different compiler adapters.

The following parent coordinate names are local to this section:
\[
\begin{gathered}
 J,F,\alpha,z,f,h,i,o,j,s,w,\tau,\eta,\zeta,y,\gamma,\\
 N,E,P,c,a,X,D,C_2,L,F_2,V,g,I,B_{\mathrm{aux}},\Delta .
\end{gathered}
\]
All 31 coordinates are strictly positive. Put
\[
 q=\beta J+1,\qquad k=\eta+\zeta,\qquad
 H_p=4a+3,\qquad C_p=q-F-\alpha.
\]
The 22 equations, in their original order, are as follows.

\begin{longtable}{@{}r@{\qquad}l@{}}
\toprule
Row & Equation\\
\midrule
\endhead
1 & \(L=q^2-1\)\\[2pt]
2 & \(N=q(L+1)\)\\[2pt]
3 & \(X=wq\)\\[2pt]
4 & \(a=E+sN\)\\[2pt]
5 & \(E=X(a-E)\)\\[2pt]
6 & \(c=k(a-E)+\eta\)\\[2pt]
7 & \(P=E(c-\eta)+k\)\\[2pt]
8 & \(\tau^2=P(P-k)+1\)\\[2pt]
9 & \(D=X+ac+\gamma H_p\)\\[2pt]
10 & \(\Delta=a^2+H_p\)\\[2pt]
11 & \(C_2=c^2\)\\[2pt]
12 & \(D^2=\Delta C_2+1\)\\[2pt]
13 & \(k=jc-V+1+hE\)\\[2pt]
14 & \(jc-V=(q-F)(N-q)-C_pL+(M_C+qM_F)J\)\\[2pt]
15 & \((K_0+w)C_p+q-F-z(q-1)=1\)\\[2pt]
16 & \(F_2=f^2\)\\[2pt]
17 & \(V=of-c\)\\[2pt]
18 & \(g=y^2-V^2\)\\[2pt]
19 & \(I=iC_2\)\\[2pt]
20 & \(B_{\mathrm{aux}}=I^2\)\\[2pt]
21 & \(B_{\mathrm{aux}}=\Delta(F_2-1)\)\\[2pt]
22 & \(y^2-B_{\mathrm{aux}}g=1\)\\
\bottomrule
\end{longtable}

Every residual is quadratic after expanding the four affine abbreviations.
The inherited parent theorem is
\begin{equation}\label{eq:parentsemantic}
 \exists\mathbf p\in\Z_{>0}^{31}\;
       \text{all 22 equations hold}
 \quad\longleftrightarrow\quad
 K215\text{ halts on blank input}
 \quad\longleftrightarrow\quad
 \neg\Con(\ZFC).
\end{equation}
Its sources are the copied parent proof, coefficient adapter, and K215
endpoint \cite{package31,k215,native}.

For each coordinate substitute \(p_i=n_i+1\).
Form left side minus right side, and reorder the rows as
\[
 1,2,3,4,6,7,\ldots,22,5.
\]
Call the resulting polynomials \(r_1,\ldots,r_{22}\in\Z[n_1,\ldots,n_{31}]\).
The last row satisfies
\begin{equation}\label{eq:lastconstant}
 r_{22}(0,\ldots,0)=1.
\end{equation}
Thus the all-zero tuple is not a common zero, and the leading coefficient
of the later reversed packing will be positive. The successor shift is
essential: the parent has positive witnesses, whereas the source coding
uses nonnegative digits.

An exact symbolic audit compares the parent circuit with an independently
retrieved sparse representation, keeping all four ports as independent
symbols. It verifies both the shift and every coefficient used below.
Each ordinary coefficient of each normalized row is a polynomial of port
degree at most two, with sum of absolute integer coefficients at most eight.
Consequently the completely fixed positive integer
\begin{equation}\label{eq:H}
 H=8(1+\beta+K_0+M_C+M_F)^2
\end{equation}
bounds the absolute values of all ordinary coefficients of all \(r_j\).

\section{Direct quadratic packing at several target positions}
\label{sec:packing}

\subsection{A small Sidon set replaces exponential positions}

A finite set is a \(B_2\) set, or Sidon set in the convention used here, if
all its unordered pair sums, including repeated summands, are distinct.
Use the following 32 positions, in increasing order:
\begin{equation}\label{eq:weights}
\begin{split}
(w_0,\ldots,w_{31})={}&(0,10,35,55,62,105,144,145,\\
 &159,168,232,245,306,323,325,353,\\
 &371,374,382,494,499,525,566,578,\\
 &610,647,681,685,741,777,783,799).
\end{split}
\end{equation}
All \(32\cdot33/2=528\) pair sums are distinct modulo \(1023\).
They are therefore distinct as integers as well. Padding a monomial of
degree less than two with zero-weight letters makes every monomial of
degree at most two correspond to exactly one unordered pair.
It follows that
\[
 \alpha\longmapsto e(\alpha)=\sum_{i=1}^{31}\alpha_iw_i
 \quad (|\alpha|\le2)
\]
is injective.

The list was obtained from the classical Bose--Chowla construction in
\(\mathbb F_{32^2}\), with an affine change modulo \(1023\).
The generation metadata use the binary polynomial \(x^{10}+x^3+1\),
multiplier 47 and shift 879. The finite pair-sum certificate alone
establishes the property required here. Thus no finite-field
implementation is a hidden assumption of the particular equation.
Bose--Chowla constructions and their affine changes are standard
\cite{bosechowla,obryant}; global priority for this style of coding is not
claimed.

Put
\begin{equation}\label{eq:targets}
\begin{gathered}
t=799,\qquad E_j=2t+(j-1)(2t+1)\quad(1\le j\le22),\\
E_{\max}=44t+21=35177,\qquad
M=E_{\max}+2t=36775.
\end{gathered}
\end{equation}
Write \(r_j=\sum_\alpha a_{j,\alpha}n^\alpha\), omitting zero terms, and
set
\begin{equation}\label{eq:weighted}
 c_{j,\alpha}=(2-|\alpha|)!\prod_{i=1}^{31}\alpha_i!\;a_{j,\alpha},
 \qquad
 \Dcode(Z)=\sum_{j=1}^{22}\sum_{|\alpha|\le2}
        c_{j,\alpha}Z^{E_j-e(\alpha)}.
\end{equation}
All exponents are nonnegative. The coefficient records number 136;
there are 161 terms when their port polynomials are expanded.
The bound \(|c_{j,\alpha}|\le2H\) follows immediately.

\begin{lemma}[Simultaneous target extraction]\label{lem:extract}
For any \(\mathbf n\in\N^{31}\), put
\[
 C_{\mathbf n}(Z)=\left(1+\sum_{i=1}^{31}n_iZ^{w_i}\right)^2.
\]
Then, for every \(j=1,\ldots,22\),
\begin{equation}\label{eq:extract}
 [Z^{E_j}]\,C_{\mathbf n}(Z)\Dcode(Z)=2r_j(\mathbf n).
\end{equation}
Every coefficient of the product has absolute value at most
\[
 2H\left(1+\sum_i n_i\right)^2.
\]
\end{lemma}

\begin{proof}
The coefficient of \(Z^{e(\alpha)}\) in \(C_{\mathbf n}\) is exactly
\[
 \frac{2!}{(2-|\alpha|)!\prod_i\alpha_i!}\,\mathbf n^\alpha
\]
by injectivity of the padded pair sums. This cancels the factorial
multiplier in \eqref{eq:weighted}.

The reversed support belonging to row \(j\) lies in
\([E_j-2t,E_j]\). These intervals are disjoint because successive targets
are \(2t+1\) apart. After multiplication by \(C_{\mathbf n}\), row \(j\)
has support in \([E_j-2t,E_j+2t]\), which contains no other target.
Thus no other row contributes to the coefficient at \(E_j\).

All coefficients of \(C_{\mathbf n}\) are nonnegative and their sum is
\((1+\sum_i n_i)^2\). Every coefficient of \(\Dcode\) has absolute value
at most \(2H\), since its row supports are disjoint. The stated product
bound follows by convolution.
\end{proof}

The product support intervals may overlap \emph{away from their targets}.
No disjointness of those larger intervals is needed. Allowing that
overlap is why spacing \(2t+1\), rather than \(4t+1\), is sufficient.

\subsection{A small example}

Take \(w=(0,1,3)\), the equations
\[
 r_1=n_1+n_2-5,\qquad r_2=6-n_1n_2,
\]
and targets \(E_1=6,E_2=13\). The reversed polynomial is
\[
 \Dcode(Z)=12Z^{13}-Z^9-10Z^6+Z^5+Z^3.
\]
Multiplication by \((1+n_1Z+n_2Z^3)^2\) produces
\[
 [Z^6]C_{\mathbf n}\Dcode=2(n_1+n_2-5),\qquad
 [Z^{13}]C_{\mathbf n}\Dcode=2(6-n_1n_2).
\]
The generic support enclosures \([0,12]\) and \([7,19]\) overlap. Nevertheless,
each tested coefficient belongs to just one row.
This example is also the independent positive and negative control used
in \file{code/check_multitarget_toy.py}.

\subsection{The structural gain for arbitrary quadratic systems}

For \(n\) variables and \(r\) quadratic rows, use a \(B_2\) set of \(n+1\)
positions including zero. Bose--Chowla sets give maximum position
\(t=O(n^2)\): take a power of two \(q\ge n+1\) with \(q<2(n+1)\), choose
\(n+1\) members of the \(q\)-element construction modulo \(q^2-1\), and
translate one to zero. The same argument uses
\[
 E_j=2t+(j-1)(2t+1),\qquad
 M=(2r+2)t+r-1=O(rn^2).
\]
This is a polynomial bound on digit-position span. It replaces the
exponential span of the old general-purpose monomial coding for this
class of systems. It is a statement about this coding step, not a
complexity bound for solving the equations.

\section{The explicit source and its two directions}
\label{sec:source}

\subsection{A fixed coefficient bound}

Set once and for all
\begin{equation}\label{eq:betasource}
 \beta_S=2^{2140030}.
\end{equation}
This is an even power of two and satisfies
\begin{equation}\label{eq:betabound}
 \beta_S>4\cdot33^2 H.
\end{equation}
For a transparent bound, each native port is less than \(2^{1070004}\).
Thus \(1+\beta+K_0+M_C+M_F<2^{1070007}\),
\(H<2^{2140017}\), and \(4\cdot33^2<2^{13}\).
The exact integer check is stronger: \(H\) has 2,140,010 bits and
\(4\cdot33^2H\) has 2,140,022 bits.
The complete bytes of \(\beta_S\) are supplied, rather than just its
residue modulo a few primes.

\subsection{Two-variable polynomial source}

From this point forward the letters refer to the finalizer, not the
positive parent coordinates. Start with integer unknowns \(f,g\) and define
\begin{equation}\label{eq:source1}
\begin{aligned}
 b&=1+3f,& \BB&=\beta_Sb^2,&
 \ell&=\sum_{i=1}^{31}\BB^{w_i},\\
 N_0&=\BB^{t+1},&
 N_1&=4\BB^{M+1},&
 N&=N_0N_1.
\end{aligned}
\end{equation}
Define the masks and packed words by
\begin{equation}\label{eq:source2}
\begin{aligned}
 M_{\mathrm c}&=N_0-1-(b-1)\ell,\\
 Q_{\mathrm c}&=(\BB-2)\sum_{j=1}^{22}\BB^{E_j},\\
 T_{\mathrm c}&=M_{\mathrm c}+Q_{\mathrm c}N_0,\\
 J_{\mathrm c}&=2(1+g)^2\Dcode(\BB)+
                      \sum_{i=0}^{M}\BB^{i+1},\\
 S_{\mathrm c}&=g+J_{\mathrm c}N_0,\\
 R_{\mathrm c}&=(S_{\mathrm c}+T_{\mathrm c}+1)N+T_{\mathrm c}+1.
\end{aligned}
\end{equation}
Finally put
\begin{equation}\label{eq:source3}
 X=(N-1)R_{\mathrm c},\qquad Y=N^2,\qquad
 \capC=b\BB^t.
\end{equation}
In particular,
\[
 N=4\BB^{37576},\qquad
 \capC=\beta_S^{799}b^{1599}
       =2^{1709883970}b^{1599}.
\]
Every expression in \eqref{eq:source1}--\eqref{eq:source3} is an integer
polynomial in \(f,g\).

\begin{proposition}[Source interface]\label{prop:source}
The following statements hold.
\begin{enumerate}[label=\textup{(\roman*)}]
\item If the \(r_j\) have a common zero in \(\N^{31}\), then for every
sufficiently large positive integer \(f\) for which \(b\) is both a
square and a power of two, there is an integer \(g\) such that
\[
 b\le g<\capC,\qquad Y\mid\binom{2X}{X}.
\]
\item If \(f\ne0\), \(b\) is a square and a power of two,
\[
 0\le g<2\capC,\qquad Y\mid\binom{2X}{X},
\]
then the \(r_j\) have a common zero in \(\N^{31}\).
\item Already under \(f\ne0\), \(b\) a square, and \(0\le g<2\capC\),
without assuming that \(b\) is a power of two, one has
\[
 0\le S_{\mathrm c},T_{\mathrm c}<N,\qquad
 3\mid X,\qquad X>3b,\qquad Y>\max\{b,256\}.
\]
\end{enumerate}
\end{proposition}

\subsection{Ranges before recovering a power-of-two base}

Suppose \(f\ne0\) and \(b\) is a square. As \(b=1+3f\), this forces
\(f>0\), \(b\ge4\), and \(\BB>2b\).
The highest coefficient of \(\Dcode\), at \(E_{\max}\), is
\(2r_{22}(0)=2\). All remaining coefficients have absolute value at
most \(2H\), so
\[
 \Dcode(\BB)\ge
 2\BB^{E_{\max}}-
 2H\frac{\BB^{E_{\max}}-1}{\BB-1}>0.
\]
The digits of \(M_{\mathrm c}\) are \(\BB-b\) at the selected positions
and \(\BB-1\) elsewhere below \(N_0\). Thus \(0\le M_{\mathrm c}<N_0\).
Similarly \(0\le Q_{\mathrm c}<\BB^{E_{\max}+1}<N_1\), and hence
\(0\le T_{\mathrm c}<N\).

For an arbitrary integer \(0\le g<2b\BB^t\), integrality gives
\[
 1+g\le2b\BB^t,\qquad g<N_0.
\]
The coefficient bound implies \(16Hb^2\le\BB-2\). Therefore
\begin{align*}
0<J_{\mathrm c}
&\le16Hb^2\BB^{2t}\sum_{i=0}^{E_{\max}}\BB^i
                      +\sum_{i=0}^{M}\BB^{i+1}\\
&<\frac{16Hb^2+\BB}{\BB-1}\BB^{M+1}
\le2\BB^{M+1}<N_1 .
\end{align*}
It follows that \(0\le S_{\mathrm c}<N\). This argument works for every
integer in the broad permitted interval, before any digit recovery.

Since \(\beta_S\equiv b\equiv1\pmod3\), we have
\(\BB\equiv1\pmod3\) and \(N=4\BB^{37576}\equiv1\pmod3\).
Thus \(3\mid X\). Further,
\(R_{\mathrm c}\ge N+1>\BB>b\) and \(N\ge4\), so \(X>3b\).
The displayed fixed base also gives \(Y=N^2>\max\{b,256\}\).
This proves part (iii).

\subsection{Recovering the digits and testing every row}

Assume additionally that \(b\) is a power of two. Then
\(\BB,N_0,N\) are powers of two. Let \(\tau_2(u,v)\) be the number of
binary carries when adding nonnegative integers \(u,v\).
Equivalently,
\[
 \tau_2(u,v)=\pc(u)+\pc(v)-\pc(u+v),
\]
where \(\pc\) denotes the number of 1-bits.
In particular, \(\tau_2(u,v)=0\) means their 1-bits are disjoint.

For \(0\le g<2\capC<N_0\), the low mask gives
\begin{equation}\label{eq:digitmask}
 \tau_2(g,M_{\mathrm c})=0
 \quad\Longleftrightarrow\quad
 g=\sum_{i=1}^{31}n_i\BB^{w_i},\qquad 0\le n_i<b.
\end{equation}
An unselected block is masked by \(\BB-1\), forcing its digit to zero.
A selected block is masked by \(\BB-b\); its available low bits represent
exactly the integers below \(b\).

For a word of this form, define
\[
 K_{\mathrm c}=(1+g)^2\Dcode(\BB)
                  +\frac{\BB}{2}\sum_{i=0}^{M}\BB^i,
 \qquad J_{\mathrm c}=2K_{\mathrm c}.
\]
This expression is used only in the proof; it is not a division gate in
the polynomial circuit. By Lemma \ref{lem:extract},
\[
 \bigl|[Z^i]C_{\mathbf n}\Dcode\bigr|
 \le2H(1+\textstyle\sum_i n_i)^2<\BB/2.
\]
Hence the base-\(\BB\) digits of \(K_{\mathrm c}\) are exactly
\[
 \BB/2+[Z^i]C_{\mathbf n}(Z)\Dcode(Z),
 \quad 0\le i\le M,
\]
all strictly between zero and \(\BB\).
In particular, its digit at \(E_j\) is \(\BB/2+2r_j(\mathbf n)\).
For \(1\le e<\BB\), the no-carry condition
\[
 \tau_2(e,\BB/2-1)=0
\]
holds exactly when \(e=\BB/2\). Applying that test at every target, and
then shifting both summands left by one bit, gives
\begin{equation}\label{eq:targettest}
 \tau_2(J_{\mathrm c},Q_{\mathrm c})=0
 \quad\Longleftrightarrow\quad
 r_1(\mathbf n)=\cdots=r_{22}(\mathbf n)=0.
\end{equation}
The lower and upper masks occupy separate blocks, so
\[
 \tau_2(S_{\mathrm c},T_{\mathrm c})=0
 \quad\Longleftrightarrow\quad
 \tau_2(g,M_{\mathrm c})=\tau_2(J_{\mathrm c},Q_{\mathrm c})=0.
\]

\subsection{An elementary central-binomial bridge}

\begin{lemma}\label{lem:binomial}
Let \(N=2^e\) and \(0\le S,T<N\). Put
\[
 X=(N-1)\bigl((S+T+1)N+T+1\bigr).
\]
Then
\[
 N^2\mid\binom{2X}{X}
 \quad\Longleftrightarrow\quad \tau_2(S,T)=0.
\]
\end{lemma}

\begin{proof}
The exact base-\(N\) decomposition is
\[
 X=(S+T)N^2+(N-1-S)N+(N-1-T).
\]
The two low blocks are \(e\)-bit complements. It follows that
\[
 \pc(X)=2e+\pc(S+T)-\pc(S)-\pc(T)
       =2e-\tau_2(S,T).
\]
Legendre's formula \(\nu_2(n!)=n-\pc(n)\) gives
\(\nu_2\binom{2X}{X}=\pc(X)\).
This valuation is at least \(2e\) exactly when the nonnegative carry
count is zero.
\end{proof}

This is the binary specialization of the carry construction used by Sun
\cite{sun2021}; the proof above supplies the exact identity needed here.
Combining the lemma with \eqref{eq:digitmask} and
\eqref{eq:targettest} proves the reverse direction of
Proposition \ref{prop:source}.

For completeness, start with a common zero \(\mathbf n\). Choose
\(b=2^{2q}>\max_i n_i\), as large as desired, and put \(f=(b-1)/3\).
This is a positive integer. Set \(g=\sum_i n_i\BB^{w_i}\).
The tuple is nonzero by \eqref{eq:lastconstant}; hence
\(g\ge\BB\ge b\). The geometric-series bound and \(\BB>b\) give
\(g<b\BB^t=\capC\). Both masks pass, so Lemma \ref{lem:binomial} yields
the required binomial divisibility. This proves part (i) and completes
the source proof.

\section{The complete ten-integer polynomial}
\label{sec:equation}

The source just proved replaces the source of package 31. All subsequent
polynomials are retained from its structural three-square finalizer
\cite{package31}, whose Lucas ingredients come from \cite{sun2021,sun2026}.
Introduce the remaining eight integer unknowns \(h,k,l,w,x,y,z,m\).
The complete list of unknowns is now exactly
\[
 f,g,h,k,l,w,x,y,z,m.
\]

Starting with \(b,X,Y,\capC\) from Section \ref{sec:source}, define
\begin{equation}\label{eq:bridge}
\begin{aligned}
 L&=lY,& U&=2LX,& V&=4gwY,& W&=bw,\\
 A&=U(V+1),& B_1&=2X+1,& C&=B_1+(A-2)h,\\
 P_n&=U^2V,& K&=X+1+k(P_n-2),\\
 S&=2A-5,& T&=3WC-2(W^2-1),\\
 R_2&=(P_n^2-4)K^2+4.
\end{aligned}
\end{equation}
Put
\begin{equation}\label{eq:positivity}
\begin{aligned}
 \Delta_b&=8\capC^3gK^2-
       g^2\bigl(32(C-KL)^2\capC^3+g^2K^2\bigr),\\
 O&=f^2l^2\Delta_b,\qquad \xi=zx.
\end{aligned}
\end{equation}
The modified Lucas certificate is
\begin{equation}\label{eq:lucas}
\begin{aligned}
 D&=(A^2-4)C^2+4,& E&=C^2D\xi,\\
 F&=4(A^2-4)E^2+1,&
 G^*&=2(A+2)(A-2)^2E^2-1,\\
 \Omega&=16E,&
 H^*&=C+(B_1-C)F+\Omega Fy,\\
 I^*&=((G^*)^2-1)(H^*)^2+1.
\end{aligned}
\end{equation}
The divisibility merger and radicands are
\begin{equation}\label{eq:merger}
\begin{aligned}
 \Sigma&=\xi S^5,& r&=T+Sy,\\
 \Theta&=(r^3-rS^2-S^3)(r^2+23S^2),\\
 r_0&=b\bigl(1+(3O-4)b^2z^2\bigr),&
 r_1&=R_2,&r_2&=DFI^*,\\
 v_1&=r_0,& v_2&=r_0(P_n^2+K^2+3).
\end{aligned}
\end{equation}
Finally set
\begin{equation}\label{eq:norminputs}
\begin{aligned}
 Z_{\mathrm n}&=\Sigma m-\Theta,\\
 a_0&=\Sigma^2r_0,&
 a_1&=\Sigma^2v_1^2r_1,&
 a_2&=\Sigma^2v_2^2r_2,\\
 V_{\mathrm n}&=Z_{\mathrm n}^2+a_2-a_0-a_1,\\
 W_{\mathrm n}&=V_{\mathrm n}^2+4Z_{\mathrm n}^2a_2-4a_0a_1.
\end{aligned}
\end{equation}
The promised equation is the following single equation:
\begin{equation}\label{eq:final}
 \boxed{\quad
 F_{10}=W_{\mathrm n}^2-16V_{\mathrm n}^2Z_{\mathrm n}^2a_2=0.
 \quad}
\end{equation}
Equations \eqref{eq:bridge}--\eqref{eq:norminputs} are substitution rules,
not additional equations with additional unknowns. In particular,
\eqref{eq:final} contains no radicals, variable exponents, divisibility
predicates, order conditions, or hidden parameters.

\section{Why the finalizer preserves the equivalence}
\label{sec:finalproof}

\subsection{The inherited bridge in the interface form used here}

Let \(u_n(A)\) be the Lucas sequence \(u_0=0,u_1=1\),
\(u_{n+1}=Au_n-u_{n-1}\), with its usual extension to integer indices.
The formulas above are those of Sun's modified binomial bridge, followed
by the package-31 changes.

The literal statements of Sun's 2026 Propositions 6.1 and 6.2 refer to
his earlier source recipe. We use their \emph{proofs parametrically}.
In those proofs the source is used only through: its forward and reverse
binomial-divisibility implications; \(b=1+3f\) (the parameter \(a=0\));
the fixed form \(\capC=2^c b^d\) with \(c=1709883970\) even and \(d=1599\)
odd; and the size and congruence conditions
in Proposition \ref{prop:source}. All these properties were proved above.
No appeal to the old quartic coding identity is left in the composition.

In this form the inherited bridge has the following consequences:
\begin{enumerate}
\item From a common nonnegative parent zero, there are
\(f,g,h,k,l,w\) with \(A>2\) even, \(B_1>0\) odd,
\[
 b\in\Sq(\Z),\quad b\le g<\capC,\quad
 C=u_{B_1}(A),\quad R_2\in\Sq(\Z),\quad S\mid T,\quad O>0.
\]
\item Conversely, for the displayed formulas, the conditions
\[
 b\in\Sq(\Z),\quad O>0,\quad \xi\ne0,\quad
 DFI^*\in\Sq(\Z),\quad R_2\in\Sq(\Z),\quad S\mid T
\]
imply a common nonnegative parent zero.
\item If \(A>2\) is even, \(B_1>0\) is odd, and \(C=u_{B_1}(A)\), then
for every positive integer \(d\) there is \(\xi>0\) divisible by \(d\)
such that \emph{each} residue class modulo \(\xi\) contains a \(y\)
making \(DFI^*\) a square.
\end{enumerate}
The last assertion is the full residue-class interpolation theorem
\cite[Theorem 4.1]{sun2026}; the first two use its Propositions 6.1--6.2
and the cited Sun-2021 Lucas estimates.

For clarity, the reverse proof first obtains \(f,l\ne0\),
\(g>0\), \(K\ne0\), and
\[
 0<g<2\capC,\qquad 4(C-KL)^2<K^2
\]
from \(O>0\). The already established source ranges then apply.
The opening argument of Sun's 2021 Theorem 4.13 gives
\(W\ne0\) and \(1<B_1<|A|/2-1\). Also \(A\) is even, \(B_1\) is odd, and
\(A-2\mid C-B_1\) by definition. These are the applicability conditions
for the reverse modified Lucas certificate. It recovers
\(C=u_{B_1}(A)\), and the remaining
Lucas and approximation arguments give
\[
 K=u_{X+1}(P_n),\qquad W=2^{B_1},\qquad
 Y\mid\binom{2X}{X}.
\]
As \(W=bw\) and \(b>0\), \(b\) is a power of two.
Proposition \ref{prop:source}(ii) finishes that implication.
This ordering is why the source bounds for arbitrary permitted \(g\)
were needed before power-of-two recovery.

\subsection{The quintic divisibility merger}

The package-31 merger uses
\[
 \pi(v)=(v^3-v-1)(v^2+23).
\]
It has no rational zero but has a root modulo every positive integer.
Here is the local argument. The cubic has discriminant \(-23\).
For an odd prime \(p\ne23\), a root of the cubic is simple. If it has
no root, it is irreducible of degree three, so Frobenius acts as an even
3-cycle and its discriminant is a square modulo \(p\). Then the
quadratic has a simple root. At \(p=23\), the cubic has the simple root
3. These roots lift to every odd prime power. At 2, \(-23\equiv1\pmod8\)
is a square in the 2-adic integers. The Chinese remainder theorem
combines the prime-power roots. The rational-root test excludes a
rational root of the cubic, and the quadratic is positive on \(\mathbb Q\).

The integer \(S=2A-5\) is always odd and nonzero.
If \(S\mid T\), write \(T=Sq\). For any \(\xi\ne0\), choose a root
\(c\) of \(\pi\) modulo \(|\xi|\) and prescribe
\[
 y\equiv c-q\pmod{\xi}.
\]
Then \(\Theta=S^5\pi(q+y)\), so \(\Sigma=\xi S^5\mid\Theta\).
Conversely, if \(S^5\mid\Theta\) but \(S\nmid T\), choose a prime
with \(\alpha=\nu_p(S)>\beta=\nu_p(T)\).
The two homogeneous factors of \(\Theta\) have valuations \(3\beta\)
and \(2\beta\); their leading terms have uniquely least valuation.
Thus \(\nu_p(\Theta)=5\beta<5\alpha\), a contradiction.
Furthermore \(\Theta\ne0\), because its vanishing would give a rational
zero \((T+Sy)/S\) of \(\pi\).

\subsection{The coprime square merger}

Suppose \(z\ne0\) and \(r_0\) is a square.
The factors \(b\) and \(1+(3O-4)b^2z^2\) are coprime.
The second factor cannot vanish: vanishing forces
\(|b|=|z|=1\), while \(b\equiv1\pmod3\) then gives \(b=1\), \(f=0\),
and \(O=0\), contradicting the vanishing equation.
Unique factorization therefore makes \(b\) a signed square.
Since a negative nonzero square cannot be 1 modulo 3, \(b\) is a
positive square and the second factor is also a square.
If \(O\le0\), that factor is at most
\(1-4(bz)^2\le-3\), impossible. Thus \(O>0\).

For the converse, use positive bridge witnesses.
The positive integer \(\Delta_b\) is a multiple of \(g\), so
\[
 O=f^2l^2\Delta_b\ge g\ge b\ge4.
\]
The integer \(d=3O-4\) is positive and is 2 modulo 3, hence nonsquare.
A positive Pell unit for \(d\), raised to a suitable finite order
modulo \(b\), gives
\[
 u^2-dv^2=1,\qquad v>0,\qquad b\mid v.
\]
Set \(z=v/b>0\). The bracket in \(r_0\) is \(u^2\), so \(r_0\) is a
square. Interpolation may now choose a positive \(\xi\) divisible by
this \(z\), followed by a \(y\) in the required residue class.
The order of these choices is legitimate; no residue class is prescribed
before its modulus exists.

\subsection{Separating three square roots}

For every integer assignment with \(z\ne0\), one has
\(|r_0|>1\). Indeed, \(r_0=0,\pm1\) would force \(|b|=1\);
then \(b=1,f=O=0\), but \(r_0=1-4z^2\le-3\).
Also
\[
 \sqrt{|r_1|}\le P_n^2+K^2+2.
\]
For example, square the right side and subtract
\((P_n^2+4)K^2+4\), an upper bound on \(|r_1|\); the remainder is
\[
 P_n^4+K^4+P_n^2K^2+4P_n^2\ge0.
\]
It follows that
\begin{equation}\label{eq:separation}
 |v_1|>\sqrt{|r_0|},\qquad
 |v_2|>\sqrt{|r_0|}+|v_1|\sqrt{|r_1|}.
\end{equation}

These inequalities imply that a rational number of the form
\[
 \varepsilon_0\sqrt{r_0}+
 \varepsilon_1v_1\sqrt{r_1}+
 \varepsilon_2v_2\sqrt{r_2},
 \qquad \varepsilon_i\in\{1,-1\},
\]
can occur only when all three radicands are integer squares.
To see this, work in the multiquadratic extension generated by the roots.
If a root is irrational, an automorphism changes it; take the
largest-index changed nonzero term. Its change has absolute value at
least twice its weight, since a nonzero integral radicand has root
absolute value at least one. The sum of the earlier changes has
strictly smaller absolute value by \eqref{eq:separation}.
Thus the entire sum changes and is not rational.
The same argument covers negative radicands, using complex absolute
value, and does not require independent square classes.
An integral radicand with rational square root is a nonnegative
integer square.

\subsection{The norm and both implications}

As an exact polynomial identity in \(Z_{\mathrm n},a_0,a_1,a_2\),
the right side of \eqref{eq:final} is
\begin{equation}\label{eq:normproduct}
 \prod_{\varepsilon\in\{1,-1\}^3}
 \left(Z_{\mathrm n}+\varepsilon_0\sqrt{a_0}
                     +\varepsilon_1\sqrt{a_1}
                     +\varepsilon_2\sqrt{a_2}\right).
\end{equation}
The radicals cancel. Repeatedly multiplying conjugate factors gives
exactly \(W_{\mathrm n}^2-16V_{\mathrm n}^2Z_{\mathrm n}^2a_2\);
the verifier also expands and compares both 35-term integer
polynomials symbolically.

Given a common nonnegative parent zero, take positive bridge witnesses.
Choose \(z>0\) by the Pell step, \(\xi>0\) divisible by \(z\) by
interpolation, and put \(x=\xi/z\). Prescribe \(y\) using the quintic
local root and then use interpolation in that residue class.
All \(r_i\) are squares, \(\Sigma\ne0\), and \(\Sigma\mid\Theta\).
Choose \(m\) so that
\[
 \Sigma m-\Theta
 =-\Sigma\bigl(\varepsilon_0\sqrt{r_0}
               +\varepsilon_1v_1\sqrt{r_1}
               +\varepsilon_2v_2\sqrt{r_2}\bigr)
\]
for any selection of signs. This is an integer choice.
One factor of \eqref{eq:normproduct} vanishes, giving \(F_{10}=0\).

Conversely suppose \(F_{10}=0\). We always have \(S\ne0\) and
\(\Theta\ne0\). If \(\Sigma=0\), all \(a_i=0\) and the polynomial
reduces to \(\Theta^8\ne0\), a contradiction. Therefore
\(\xi=zx\ne0\) and \(z\ne0\).
Dividing a vanishing norm factor by \(\Sigma\) makes a weighted signed
sum of the roots of the \(r_i\) rational. Separation forces all three
radicands to be squares, and the sum is then an integer. Hence
\(\Sigma\mid\Theta\). The quintic merger gives \(S\mid T\), and the
coprime square merger gives \(b\) square and \(O>0\).
The reverse bridge now supplies a common nonnegative parent zero.
Combining with \eqref{eq:parentsemantic} proves \eqref{eq:main}.
Classical negation gives the assertion about absence of solutions.

\section{Exact degree and the 782-operation circuit}
\label{sec:degree}

\subsection{Evaluation without expanding the enormous polynomial}

The literal polynomial is specified by
\file{data/polynomial.json}, a finite directed acyclic graph whose
non-leaf gates are binary addition, subtraction, and multiplication.
The fixed nonliteral leaves have exact finite recipes in
\file{data/coefficient_recipes.json}. There are 16 such leaves,
including \(\beta_S\); the ordinary small integer leaves are literal.
The four native ports are not unknowns.

The builder evaluates \(\Dcode\) efficiently by homogenizing each
quadratic row. Put
\[
 U_0=\BB^t,\qquad U_i=\BB^{t-w_i},\qquad
 Q_j=\sum_\alpha c_{j,\alpha}
                   U_0^{2-|\alpha|}\prod_i U_i^{\alpha_i}.
\]
Then the exact identity
\[
 \Dcode(\BB)=\sum_{j=1}^{22}
       \BB^{(j-1)(2t+1)}Q_j
\]
allows shared multivariate Horner evaluation of the row forms and
ordinary Horner evaluation across rows. It preserves the literal
coefficient-extraction polynomial.

Fixed powers use shared repeated squaring. For a geometric sum
\(G_n=1+\BB+\cdots+\BB^{n-1}\), the builder uses
\[
 G_{2n}=G_n(1+\BB^n),\qquad
 G_{2n+1}=G_{2n}+\BB^{2n}.
\]
Thus the large sum limits do not entail one gate per summand.
Every fixed power is expanded into binary gates before counting.

The operation convention matches package 31:
each binary addition, subtraction, or multiplication costs one;
fixed leaves, arithmetic involving only fixed constants, and reuse
are free. Multiplication of a variable-dependent expression by a
fixed numeral is charged. Under this convention the live graph has
\[
 221+25+536=782
\]
operations. This is a certified count of the supplied evaluator, not
an arithmetic-circuit lower bound.

\subsection{A symbolic proof of exact degree}

Let
\[
 u=M+t+2=47t+23=37576.
\]
Every fixed coefficient has degree zero and every unknown has degree one.
The highest term of \(\Dcode(\BB)\) is nonzero because the last normalized
row has constant coefficient one. The geometric sum in \(J_{\mathrm c}\)
has degree \(2(M+1)\), strictly greater than the other summand's degree.
It follows successively that
\[
\begin{array}{c|cccccc}
\text{expression} &N&S_{\mathrm c}&T_{\mathrm c}&X&Y&\capC\\
\hline
\text{degree}&2u&2u&2u-4t&6u&4u&2t+1 .
\end{array}
\]
In particular,
\[
 \deg X=225456,\qquad\deg Y=150304,\qquad
 \deg\capC=1599.
\]

Write \(d_X=\deg X\), \(d_Y=\deg Y\), and \(d_C=\deg\capC\).
Direct propagation through the retained finalizer gives
\begin{align*}
 \deg Z_{\mathrm n}&=5d_X+10d_Y+30,\\
 \deg a_0&=14d_X+28d_Y+3d_C+57,\\
 \deg a_1&=26d_X+48d_Y+6d_C+98,\\
 \deg a_2&=114d_X+224d_Y+6d_C+440.
\end{align*}
The last degree strictly exceeds the degrees of \(a_0,a_1,Z_{\mathrm n}^2\).
Consequently the highest homogeneous part of \(F_{10}\) is the fourth
power of the highest homogeneous part of \(a_2\), which is nonzero.
Hence
\begin{align}
 \deg F_{10}
 &=456d_X+896d_Y+24d_C+1760 \notag\\
 &=6320u+48t+1784 \notag\\
 &=297088t+147144
 =237520456. \label{eq:degree}
\end{align}
Appendix \ref{app:degree} records the intermediate degrees and the
strict-dominance argument for the few expressions with subtraction.

There is also an independent exact lower-bound certificate. After
substituting all ten variables by one indeterminate \(T\), the
coefficient at the circuit's total-degree upper bound is
\[
 [T^{237520456}]F_{10}(T,\ldots,T)
 \equiv449536588\pmod{1000000007}.
\]
This nonzero residue proves attainment of the upper bound. It is a
coefficient computation, not random sampling or numerical interpolation.

\subsection{Coefficient size is a separate measure}

The new source constant has 2,140,031 bits, whereas package 31's old
source constant \(2^{4269882}\) has 4,269,883 bits.
The four native compiler ports are exactly unchanged. These numbers
describe fixed leaves in the circuit representation. They are
\emph{not} the heights of coefficients after fully expanding \(F_{10}\).
Those expanded heights have not been measured, and no claim about their
optimality is made.

\section{Verification, reproducibility, and proof status}
\label{sec:verification}

From the archive root, run
\begin{verbatim}
python3 code/verify.py
pdflatex -interaction=nonstopmode -halt-on-error report.tex
pdflatex -interaction=nonstopmode -halt-on-error report.tex
\end{verbatim}
The mathematical build and audits use only the Python standard library.
No network access or proof search is required. Run Python without
\file{-O}, since assertions are part of the checking process.

The verifier reconstructs the coefficient payloads, regenerates the
polynomial, and checks the following distinct claims.

\begin{longtable}{@{}p{0.31\linewidth}p{0.63\linewidth}@{}}
\toprule
Check & What the result establishes\\
\midrule
\endhead
Native tuple binding &
All four raw payload hashes equal package 31's pinned tuple; 430 literal
machine actions agree with the transition JSON.\\[4pt]
Native finite geometry &
69,230 exact scalar identities from the inherited coefficient builder
are rechecked. They certify its finite geometry, not the entire
unbounded compiler theorem.\\[4pt]
Parent and target algebra &
Exact sparse polynomial comparison with independent parent data; all
528 pair sums; all 22 target-coefficient identities; normalized row and
coefficient bounds.\\[4pt]
Independent circuit evaluator &
No imports from the builder; 34,848 modular coefficient diagnostics and
6,720 named-node comparisons in 120 complete circuit cases. These are
transcription checks, not universal polynomial identity proofs.\\[4pt]
Exact norm identity &
All eight signed factors expand to the stated 35-term radical-free
integer polynomial.\\[4pt]
Exact degree &
Structural leading-term proof plus nonzero diagonal leading
coefficient modulo a prime.\\[4pt]
Small source controls &
672 exact cases, four accepting packed tuples across the two chosen
bases, and 400 rejected malformed words. The tests use a small soluble
two-equation system.\\
\bottomrule
\end{longtable}

All checks completed successfully. The machine is not run until it
halts, and none of these programs decides whether a ZFC refutation
exists. They are practical evidence against coding or transcription
errors around a mathematical reduction with an unbounded existential
meaning.

The core files are:
\begin{itemize}
\item \file{report.tex}, \file{report.pdf}: this proof and its rendered form;
\item \file{data/polynomial.json}: the literal ten-variable equation;
\item \file{data/weights.json}, \file{data/coefficient_recipes.json}:
the finite source data and exact fixed-leaf recipes;
\item \file{coefficients/}: complete native tuple, builder, proof adapter,
literal tile and transition tables, and upstream notices;
\item \file{code/}: builder, exact algebraic checks, independent evaluator,
and one-command verification;
\item \file{certificates/}: results, hashes, operation and degree receipts;
\item \file{upstream/}: the pinned parent data, inherited finalizer proof,
K215 endpoint and literal machine, and focused mathematical audit notes.
\end{itemize}
The main result uses no artifact from \file{research_candidate/}.
That directory has a different, explicitly incomplete role described next.

\section{Further reductions worth pursuing}
\label{sec:future}

\subsection{Reduce the weight span or exploit the actual supports}

For this finalizer and 22 rows, equation \eqref{eq:degree} shows that
lowering \(t\) by one lowers the degree by exactly 297,088.
The present \(t=799\) is a verified construction, not an optimum
32-mark ruler. A shorter certified \(B_2\) set immediately improves the
same proof and need not change any inherited theorem.

More substantially, a full \(B_2\) set separates \emph{all} quadratic
monomials, while only 136 weighted records occur in this parent.
A support-aware construction might permit some collisions provided
they cannot contaminate any of the 22 required targets. Such a proposal
must replace the full target-extraction proof, rather than merely
demonstrate several successful evaluations. Jointly choosing weights,
row order, and nonuniform target gaps is a concrete finite optimization
problem with exact symbolic certificates available.

\subsection{Reuse the new source on other finalizers}

The generic source proposition is separable from the Sun-2026
ten-integer route. Rebuilding the published Sun-2021 route or the
repository's other low-variable finalizers with a direct system source
may reduce their degrees substantially. These are proposed applications,
not additional certified results of this archive. Likewise, eliminating
another positive parent witness would require a complete sign and
reconstruction argument; no such elimination is claimed here.

\subsection{A smaller native coefficient candidate, with an open obligation}

A parallel line of work factors the receiver part of the tile compiler.
It replaces 488 receiver selectors by two mode selectors and 17 Boolean
payload selectors. The candidate has 457 selectors in total and admits
materialized coefficient payloads with radix exponent 143 and cell
length 4829. Its four raw payloads total 346,098 bytes, compared with
506,664 bytes in package 36, a \(31.69\%\) decrease.

Those numbers describe a candidate, not an established replacement tuple.
The local decoder and finite convolution checks pass, but a new empty
cell may contain nonzero payload bits while its occupancy bits vanish.
The inherited binary-alignment and upper-dummy argument must be audited
or extended to cover those patterns. Without that unbounded soundness
step, the candidate cannot be substituted into the consistency equation.
The exact outstanding obligation, proposed decoder, code, notices, and
finite receipts are retained in
\file{research_candidate/FACTORIZED_RECEIVERS.md}.

Keeping this obligation visible is useful: a smaller finite coefficient
file alone does not establish the compiler's semantic equivalence.
The main construction avoids that issue by using the unchanged,
already documented package-31 tuple.

\subsection{Formalize the new interface}

The natural next proof-assistant target is the generic simultaneous-target
lemma and source proposition, with a finite \(B_2\) certificate as input.
The exact parent polynomials and norm identity then become finite
reflection checks. A full end-to-end formalization would additionally
need the native compiler and the Lucas/interpolation bridges, whose
ordinary mathematical statements are still explicit dependencies.

The present result is a substantial structural and quantitative
improvement over the reviewed ten-integer construction. It does not
prove the shortest possible equation, a lower bound on unknowns,
or worldwide priority for digit coding of systems. Classical
Sidon constructions, Sun's coding methods, and prior work on
existential arithmetization of systems remain relevant antecedents
\cite{bosechowla,obryant,sun2021,matiyasevich2009}.

\appendix
\section{The four fixed native coefficient ports}
\label{app:constants}

This appendix restates the literal arithmetic recipe of the retained
package-31 coefficient adapter \cite{package31}. It is included to make
clear that the coefficients are finite specified integers, not parameters
whose existence depends on an undecided statement.

Put
\[
 \rho=2^{121},\qquad L_{\mathrm{cell}}=8801,\qquad
 d=1064921,\qquad B_{\mathrm{native}}=2^d.
\]
The ordered 943 tiles are the literal quadruples
\((x_{i0},x_{i1},x_{i2},x_{i3})\), in west, north, east, south order,
in \file{coefficients/upstream/wang943_tiles.json}.
The included builder reconstructs this list from the 430 literal K215
actions in \file{coefficients/upstream/k215_transitions.json}.
Selector 0 is the start tile, selector 1 the impossible tile, and the
remaining 941 tiles are lexicographically ordered.

Define the 989-element support
\[
\begin{split}
\mathcal E={}&\{0,\ldots,942\}\cup\{983,\ldots,992\}
 \cup\{1003,\ldots,1032\}\\
 &{}\cup\{1063,1064,1095,1096,2195,3296\}.
\end{split}
\]
Let \(W_j=2^{10+2j}\) for \(0\le j<44\). For a tile selector \(i\), set
\[
 c_i=1+\sum_{t=0}^{3}\sum_{j=0}^{9}
       \operatorname{bit}_j(x_{it}+512)\,W_{10t+j}
       +W_{40}+W_{41}+W_{42}+W_{43}.
\]
At copy positions \(\operatorname{offset}_t+j\), for
\[
 (\operatorname{offset}_0,\ldots,\operatorname{offset}_3)
   =(1023,1013,1003,983),
\]
put \(c_{\operatorname{offset}_t+j}=W_{10t+j}\).
The four anchor coefficients are
\[
 (c_{1095},c_{1096},c_{2195},c_{3296})
       =(W_{40},W_{41},W_{42},W_{43}).
\]
Finally put \(c_{1063}=c_{1064}=0\).

Define
\[
 \mu=1022+\sum_{j=0}^{43}W_j,\qquad
 Z_p=(2^{119}-1)\sum_{e=8792}^{8798}\rho^e
                       +(2^{27}-1)\rho^{8799},
\]
and
\begin{align*}
 D_C&=\rho^{30}+\rho^{62}+\rho^{1100}+\rho^{2200}
       +\sum_{e\in\mathcal E}c_e
                    (\rho^{4393-e}+\rho^{5495-e}),\\
 D_R&=\rho^{42},\\
 M_C&=B_{\mathrm{native}}-1
       -\sum_{\substack{e\in\mathcal E\\e\ne1}}\rho^e
       -2\rho^{1063},\\
 M_F^{\mathrm{native}}
    &=4+\mu(\rho^{4393}+\rho^{5495})
       +\sum_{j=0}^{9}(\rho^{1013+j}+\rho^{1065+j})
       +\rho^{2195}+\rho^{3296}+Z_p.
\end{align*}
The four ports, as one indivisible tuple, are
\[
 \boxed{\quad
 \beta=B_{\mathrm{native}}-1,\qquad
 K_0=D_C+B_{\mathrm{native}}D_R,\qquad
 M_C,\qquad
 M_F=M_F^{\mathrm{native}}+B_{\mathrm{native}}-1.
 \quad}
\]
The materialized payloads are unsigned little-endian integers;
the \file{.bin.gz} files decompress to those exact bytes.
The coefficient builder and binding manifests specify their hashes.
The literal machine file has SHA-256
\begin{center}\small\ttfamily
fdbcef0825886035c3091c8537a69535bc96f9e265f0ae45863543d40ae70090.
\end{center}
The newly introduced source constant \(\beta_S\), also unsigned
little-endian, has SHA-256
\begin{center}\small\ttfamily
b5505995484d5993ce1c0bba8f6fec9d8a38e2d8c5b7b2ac535557bcf48dfdfe.
\end{center}
Checksums bind files; the arithmetic recipe and inherited compiler proof
give their mathematical meaning.

\section{Detailed degree ledger}
\label{app:degree}

Let \(x=d_X\), \(y=d_Y\), \(c=d_C\) in this appendix only. These letters
denote degrees, not the polynomial unknowns.

\begin{longtable}{@{}ll@{}}
\toprule
Expression & Exact degree\\
\midrule
\endhead
\(L\) & \(y+1\)\\
\(U\) & \(x+y+1\)\\
\(V\) & \(y+2\)\\
\(W\) & \(2\)\\
\(A\) & \(x+2y+3\)\\
\(B_1\) & \(x\)\\
\(C\) & \(x+2y+4\)\\
\(P_n\) & \(2x+3y+4\)\\
\(K\) & \(2x+3y+5\)\\
\(S\) & \(x+2y+3\)\\
\(T\) & \(x+2y+6\)\\
\(R_2\) & \(8x+12y+18\)\\
\(\Delta_b\) & \(4x+8y+3c+14\)\\
\(O\) & \(4x+8y+3c+18\)\\
\(r_0\) & \(4x+8y+3c+23\)\\
\(D\) & \(4x+8y+14\)\\
\(E\) & \(6x+12y+24\)\\
\(F\) & \(14x+28y+54\)\\
\(G^*\) & \(15x+30y+57\)\\
\(\Omega\) & \(6x+12y+24\)\\
\(H^*\) & \(20x+40y+79\)\\
\(I^*\) & \(70x+140y+272\)\\
\(r_1\) & \(8x+12y+18\)\\
\(r_2\) & \(88x+176y+340\)\\
\(\Sigma\) & \(5x+10y+17\)\\
\(\Theta\) & \(5x+10y+30\)\\
\(Z_{\mathrm n}\) & \(5x+10y+30\)\\
\(a_0\) & \(14x+28y+3c+57\)\\
\(a_1\) & \(26x+48y+6c+98\)\\
\(a_2\) & \(114x+224y+6c+440\)\\
\bottomrule
\end{longtable}

In \(\Delta_b\), the term of greatest degree is
\(-32g^2\capC^3K^2L^2\); its degree strictly exceeds the other terms.
In \(r=T+Sy\), the greatest-degree term comes from \(T\), so the
greatest-degree term of \(\Theta\) is the fifth power of that term.
The term \(\Theta\) strictly dominates \(\Sigma m\).
Finally \(a_2\) strictly dominates \(a_0,a_1,Z_{\mathrm n}^2\);
the greatest-degree term of the final norm is \(a_2^4\).
These observations justify equality, rather than just an upper bound,
through all potentially cancelling subtractions.

As a consistency check on the retained finalizer algebra, inserting the
old source degrees
\[
 x=120t+24,\qquad y=80t+16,\qquad c=4t+1
\]
into \(456x+896y+24c+1760\) gives
\(126496t+27064\), exactly package 31's degree formula.

\section{Attribution and integration notes}
\label{app:integration}

The new contribution is the direct quadratic-system source, its explicit
Sidon instance, the tightened multiple-target support argument, the
complete new circuit and degree calculation, and the accompanying checks.
The K215 machine, native compiler, parent system and finalizer were
obtained from ProveIt at the stated pin. The coefficient builder and DAG
utility retain their upstream provenance. Sun's variable-count theorem,
Lucas identities, binomial bridge, and residue interpolation are explicitly
inherited. Classical Sidon constructions supply the weight design.

This is an AI-assisted research package, not a claim of endorsement by
the cited authors or by the repository maintainers. Original licenses
and machine notices are included. The proposed directory name
\file{38-quadratic-packing} is an integration suggestion only; no repository
commit, publication, or upstream modification was made.

For review, the shortest useful path is:
read Lemma \ref{lem:extract}; inspect the arbitrary-\(g\) bounds in
Section \ref{sec:source}; check that the parametric bridge uses precisely
Proposition \ref{prop:source}; and run the supplied verifier.
The separate receiver candidate should remain marked incomplete until
its stated alignment obligation is proved.

\begin{thebibliography}{99}

\bibitem{repo}
V. Reshetnikov and contributors,
\emph{ProveIt: Con(ZFC) in Diophantine form}, repository collection,
revision \texttt{22e75dd6b87b1b199bb08c778d7cf075ae5aae08},
23 October 2026.
\href{https://github.com/VladimirReshetnikov/ProveIt/tree/22e75dd6b87b1b199bb08c778d7cf075ae5aae08/SetTheory/Cardinals/docs/reports/hilbert-tenth-problem/con-zfc-diophantine-forms}{Pinned collection}.
The current comparison README is copied in this archive.

\bibitem{package31}
ProveIt research package 31,
\emph{Smaller explicit Diophantine forms for Con(ZFC)},
\file{31-sun10}, at the revision in \cite{repo}.
\href{https://github.com/VladimirReshetnikov/ProveIt/tree/22e75dd6b87b1b199bb08c778d7cf075ae5aae08/SetTheory/Cardinals/docs/reports/hilbert-tenth-problem/con-zfc-diophantine-forms/31-sun10}{Pinned package}.
The inherited parent, coefficient adapter and finalizer proofs are
included with their provenance.

\bibitem{k215}
ProveIt,
\emph{K215 blank-halting and ZFC-refutation endpoint},
\file{KEndpoint215.lean}, at the revision in \cite{repo}.
\href{https://github.com/VladimirReshetnikov/ProveIt/blob/22e75dd6b87b1b199bb08c778d7cf075ae5aae08/Computability/BusyBeaver/Lean/BusyBeaver/ZFC/KEndpoint215.lean}{Endpoint source}.
The literal machine and endpoint source are included.

\bibitem{native}
ProveIt,
\emph{Native-window/Pell compiler research},
\file{native-stream-queue},
at the revision in \cite{repo}.
\href{https://github.com/VladimirReshetnikov/ProveIt/tree/22e75dd6b87b1b199bb08c778d7cf075ae5aae08/Computability/HilbertTenthProblem/Papers/research-wip/native-stream-queue}{Compiler sources}.
Used through the retained package-31 coefficient adapter and parent theorem.

\bibitem{sun2021}
Z.-W. Sun,
\emph{Further results on Hilbert's Tenth Problem},
Science China Mathematics \textbf{64} (2021), 281--306.
\href{https://arxiv.org/html/1704.03504v7}{arXiv:1704.03504v7}.
Relevant inputs: Lemma 2.6, Theorem 3.1, and Theorems 4.12--4.13.

\bibitem{sun2026}
Z.-W. Sun,
\emph{Ten unknowns for Hilbert's tenth problem over the integers},
preprint, 2 September 2026.
\href{https://arxiv.org/html/2609.01594v2}{arXiv:2609.01594v2}.
Relevant inputs: Theorems 3.1 and 4.1, Propositions 6.1--6.2.

\bibitem{bosechowla}
R. C. Bose and S. Chowla,
\emph{Theorems in the additive theory of numbers},
technical report, 1960.
\href{https://www.cs.umd.edu/~gasarch/COURSES/858/S13/BoseChowla.pdf}{Original report scan}.

\bibitem{obryant}
K. O'Bryant,
\emph{Constructing Thick \(B_h\)-Sets},
Journal of Integer Sequences \textbf{27} (2024), Article 24.1.2.
\href{https://cs.uwaterloo.ca/journals/JIS/VOL27/OBryant/obry7.pdf}{Author's article}.
See Theorem 3 and the discussion of affine changes.

\bibitem{matiyasevich2009}
Yu. Matiyasevich,
\emph{Existential arithmetization of Diophantine equations},
Annals of Pure and Applied Logic \textbf{157} (2009), 225--233.
\href{https://doi.org/10.1016/j.apal.2008.09.009}{doi:10.1016/j.apal.2008.09.009}.
Related work on coding equations and systems; no theorem from this paper
is required for the proof above.

\end{thebibliography}
\end{document}

